\section{模型评价}

在监督学习中，模型评价是衡量模型预测能力、比较不同模型并指导模型选择的重要环节。

\subsection{分类任务}
分类问题可以根据类别数划分为二分类问题与多分类问题，二分类问题即$K=2$的分类问题，多分类问题对应$K\geqslant3$的分类问题。
\subsubsection{二分类}
在二分类问题中，通常将研究中重点关注、希望模型识别出的类别称为\textbf{正类}，将另一个类别称为\textbf{负类}。\par
\begin{definition}
	\gls{ConfusionMatrix}按真实标签与预测标签的组合统计样本个数，其每个元素含义如下：
	\begin{itemize}
		\item \gls{TP}：真实为正类且预测为正类的样本数；
		\item \gls{TN}：真实为负类且预测为负类的样本数；
		\item \gls{FP}：真实为负类但预测为正类的样本数；
		\item \gls{FN}：真实为正类但预测为负类的样本数。
	\end{itemize}
\end{definition}
\begin{table}[H]
	\centering
	\caption{二分类问题的混淆矩阵}
	\begin{tabular}{ccc}
		\toprule
		真实类别$\backslash$预测类别
		&正类($\hat{y}=1$)
		&负类($\hat{y}=0$) \\
		\midrule
		正类($y=1$)
		&TP
		&FN \\
		负类($y=0$)
		&FP
		&TN \\
		\bottomrule
	\end{tabular}
\end{table}
\begin{definition}
  \gls{Accuracy}度量模型预测正确的比例：
  \begin{equation*}
    \operatorname{Accuracy}=\frac{\operatorname{TP}+\operatorname{TN}}{\operatorname{TP}+\operatorname{TN}+\operatorname{FP}+\operatorname{FN}}
  \end{equation*}
\end{definition}
Accuracy 将正类和负类的正确预测汇总为一个整体比例，因而默认两类样本具有相同的重要性，并且不同类型的分类错误具有相同的代价。然而在许多实际问题中，假正类和假负类造成的损失并不相同，我们也可能更关注模型在某一特定类别上的预测表现。\par
例如，在贷款违约预测中，若将违约客户定义为正类，则将实际违约的客户预测为不违约会使银行承担贷款损失，而将实际不违约的客户预测为违约主要意味着失去一次潜在的放贷机会。这两类错误的经济代价并不相同。因此，即使两个模型具有相同的 Accuracy，它们识别违约客户的能力也可能存在显著差异。为了考察模型针对特定类别的表现，需要进一步引入别的指标。
\begin{definition}
  \gls{Precision}度量被预测为正类的样本中有多少是真正的正类：
  \begin{equation*}
    \operatorname{Precision}=\frac{\operatorname{TP}}{\operatorname{TP}+\operatorname{FP}}
  \end{equation*}
\end{definition}
\begin{note}
  精确率关心正类预测结果的准确性，在系统对错误正类具有高成本的任务（如垃圾邮件过滤）中尤为重要。较高的精确率意味着模型产生的误报较少。
\end{note}
\begin{definition}
  \gls{Recall}又称\gls{Sensitivity}，衡量真实正类中有多少被模型成功识别，定义为：
  \begin{equation*}
    \operatorname{Recall}=\frac{\operatorname{TP}}{\operatorname{TP}+\operatorname{FN}}
  \end{equation*}
\end{definition}
\begin{note}
  召回率关注对正类的漏判率，适用于需要尽可能检出正类的场景（如疾病筛查或欺诈检测）。高召回率意味着漏报较少，但可能伴随较多假阳性，需要结合精确率来进行分析。
\end{note}
\begin{definition}
  \gls{F1Score}是精确率与召回率的调和平均，用于在二者之间取得平衡：
  \begin{equation*}
    \operatorname{F1}=2\times\frac{\operatorname{Precision}\times\operatorname{Recall}}{\operatorname{Precision}+\operatorname{Recall}}
  \end{equation*}
\end{definition}
\begin{note}
  调和平均的性质使$\operatorname{F1}$得分对低值敏感——当精确率或召回率其中之一很低时，$\operatorname{F1}$会显著下降，因此它常用作综合评价指标，适用于精确率和召回率同等重要的情形。
\end{note}
\begin{definition}
	$\operatorname{F}_{\beta}$分数是$\operatorname{F1}$得分的推广，通过参数$\beta>0$调节精确率与召回率的重要性：
	\begin{equation*}
		\operatorname{F}_{\beta}=(1+\beta^2)\frac{\operatorname{Precision}\times\operatorname{Recall}}{\beta^2\operatorname{Precision}+\operatorname{Recall}}=\dfrac{1}{\dfrac{\beta^2}{1+\beta^2}\dfrac{1}{\operatorname{Recall}}+\dfrac{1}{1+\beta^2}\dfrac{1}{\operatorname{Precision}}}
	\end{equation*}
\end{definition}
\begin{note}
	当$|\beta|=1$时，$\operatorname{F}_{\beta}$退化为$\operatorname{F1}$；当$|\beta|>1$时，召回率具有更大的权重，适用于更关注漏检的任务；当$0<|\beta|<1$时，精确率具有更大的权重，适用于更关注误报的任务。因此，相比固定地给予精确率与召回率相同权重的$\operatorname{F1}$，$\operatorname{F}_{\beta}$能够根据具体任务中两类错误的重要程度进行调整。
\end{note}
许多分类器并不直接输出类别，而是为每个样本输出一个连续得分或属于正类的概率，并在给定阈值后将其转化为分类结果。然而，分类阈值并不存在适用于所有任务的固定选择，即使模型输出的是概率，也不能机械地将$0.5$作为阈值。阈值的选择取决于正类的基础发生率、假正类与假负类的相对代价以及模型所服务的具体决策目标。\par
以罕见病筛查为例，假设某种疾病在人群中的患病率仅为$0.1\%$。若模型预测某位受检者的患病概率为$10\%$，那么其患病风险已经达到普通人群的$100$ 倍，显然应当被识别为高风险个体并接受进一步检查。然而，$10\%$仍然低于$50\%$；若机械地将$0.5$作为分类阈值，该受检者就会被判定为阴性。类似地，在身份认证或银行人脸识别中，系统可能更重视防止冒用身份，因此会采用更加严格的判定阈值，以降低错误接受非本人的概率。\par
由此可见，同一个模型在不同应用场景中可能需要采用完全不同的阈值，而模型在某一个固定阈值下的表现不能完整反映其分类能力。为了减少评价结果对特定阈值选择的依赖，可以考察阈值连续变化时模型性能的变化，并将其表示为性能曲线。我们希望一个模型在较广的阈值范围内均能保持较好的正负类区分能力，因此可以进一步使用曲线下面积评价模型的综合表现。
\begin{definition}
	称：
	 \begin{equation*}
		\operatorname{TPR}=\frac{\operatorname{TP}}{\operatorname{TP}+\operatorname{FN}},\quad
		\operatorname{FPR}=\frac{\operatorname{FP}}{\operatorname{FP}+\operatorname{TN}}
	\end{equation*}
	分别为\gls{TPR}与\gls{FPR}。
\end{definition}
\begin{definition}
  设分类器将得分不低于阈值$t$的样本预测为正类。\gls{ROC}为在阈值取值范围内连续变化$t$，对每个$t$计算对应的：
	\begin{equation*}
		\operatorname{TPR}(t)=\frac{\operatorname{TP}(t)}{\operatorname{TP}(t)+\operatorname{FN}(t)},\quad
		\operatorname{FPR}(t)=	\frac{\operatorname{FP}(t)}{\operatorname{FP}(t)+\operatorname{TN}(t)}
	\end{equation*}
	然后以$\mathrm{FPR}$为横轴、$\mathrm{TPR}$为纵轴作图得到的曲线。\gls{AUROC}定义为ROC曲线与$\mathrm{FPR}$轴和$\mathrm{FPR}=1$围成区域的面积。
\end{definition}
\begin{note}
  ROC曲线反映了阈值变化对$\operatorname{TPR}$与$\operatorname{FPR}$的权衡。$\operatorname{AUROC}$独立于具体阈值，是比较不同模型总体区分能力的常用指标，取值范围通常在$[0.5,1]$，$0.5$表示随机猜测（对任意的$t\in[0,1]$有$\operatorname{FPR}(t)=\operatorname{TPR}(t)$，模型判断正类为正类的概率和模型判断负类为正类的概率一致），$1$表示完美分类。
  \begin{figure}[H]
  	\centering
  	\begin{tikzpicture}[scale=3]
  		% 坐标轴
  		\draw[->] (0,0) -- (1.05,0) node[right] {\scriptsize FPR};
  		\draw[->] (0,0) -- (0,1.05) node[above] {\scriptsize TPR};
  		% 随机分类器的对角线
  		\draw[color=gray, dashed] (0,0) -- (1,1);
  		% 绘制示意 ROC 曲线
  		\draw[color=blue, thick] (0,0) .. controls (0.15,0.5) and (0.3,0.75) .. (1,1);
  		% 填充面积
  		\path[blue!20] (0,0) -- (0.15,0.5) .. controls (0.15,0.5) and (0.3,0.75) .. (1,1) -- (1,0) -- cycle;
  	\end{tikzpicture}
  	\caption{ROC曲线示意}
  \end{figure}
\end{note}
\begin{definition}
 	设模型将得分不低于阈值 $t$ 的样本预测为正类。\gls{PRC}为在阈值取值范围内内连续变化$t$，对每个$t$计算对应的$\operatorname{Precision}(t)$和$\operatorname{Recall}(t)$，然后以$\mathrm{Recall}$为横轴、$\mathrm{Precision}$为纵轴作图得到的曲线。\gls{AUPRC}定义为PRC曲线与$\operatorname{Recall}$轴之间区域的面积。
\end{definition}
\begin{note}
  相比$\operatorname{AUROC}$，$\operatorname{AUPRC}$更关注正类的表现。较高的$\operatorname{AUPRC}$表明模型性能较好。
  \begin{figure}[H]
  	\centering
  	\begin{tikzpicture}[scale=3]
  		% 坐标轴
  		\draw[->] (0,0) -- (1.05,0) node[right] {\scriptsize Recall};
  		\draw[->] (0,0) -- (0,1.05) node[above] {\scriptsize Precision};
  		% 绘制示意 PR 曲线
  		\draw[color=blue, thick] (0,1) .. controls (0.3,0.9) and (0.6,0.6) .. (1,0.2);
  		% 填充面积
  		\path[blue!20] (0,1) .. controls (0.3,0.9) and (0.6,0.6) .. (1,0.2) -- (1,0) -- (0,0) -- cycle;
  	\end{tikzpicture}
  	\caption{PR曲线示意}
  \end{figure}
\end{note}
对于能够输出预测概率的分类模型，模型性能至少涉及两个不同方面：一方面，模型应具有良好的排序能力，即正类样本通常获得高于负类样本的预测分数；另一方面，模型给出的预测概率还应具有可靠的概率解释。前者通常可以通过AUROC评价，但良好的排序能力并不意味着预测概率本身准确。例如，某模型可能始终给真正的正类样本赋予高于负类样本的预测概率，从而获得较高的AUROC，但对于所有预测概率约为$0.8$的样本，实际只有约$60\%$属于正类。此时模型虽然能够较好地区分样本的相对风险高低，却系统性高估了事件发生的绝对概率。因此，还需要进一步考察预测概率与实际事件发生频率之间的一致程度。
\begin{definition}
	设分类模型对样本$x_i$给出的正类预测概率为$\hat{p}_i$。将预测概率的取值范围划分为若干区间$\seq{B}{K}$，对于每个区间$B_k$，分别计算其中样本被预测为正类的概率的平均值以及真实正类比例：
	\begin{equation*}
		\operatorname{Pred}(B_k)=\frac{1}{|B_k|}\sum_{i:\hat{p}_i\in B_k}\hat{p}_i,\quad\operatorname{Obs}(B_k)=\frac{1}{|B_k|}\sum_{i:\hat{p}_i\in B_k}y_i
	\end{equation*}
	其中$|B_k|$表示落入第$k$个概率区间$B_k$中的样本数量。\gls{CalibrationCurve}为以$\operatorname{Pred}(B_k)$为横轴、$\operatorname{Obs}(B_k)$为纵轴作图得到的曲线。
\end{definition}
\begin{note}
	理想情况下校准曲线应与直线$y=x$重合。若校准曲线位于$y=x$下方，则说明模型给出的预测概率整体偏高，反之模型给出的预测概率整体偏低。
	\begin{figure}[H]
		\centering
		\begin{tikzpicture}[scale=3]
			% 坐标轴
			\draw[->] (0,0) -- (1.05,0) node[right] {\scriptsize Predicted Probability};
			\draw[->] (0,0) -- (0,1.05) node[above] {\scriptsize Observed Frequency};
			
			% 完美校准线
			\draw[dashed, gray] (0,0) -- (1,1);
			
			% 示意校准曲线
			\draw[color=blue, thick]
			(0.05,0.08)
			.. controls (0.25,0.18) and (0.40,0.32) ..
			(0.55,0.45)
			.. controls (0.70,0.55) and (0.85,0.70) ..
			(0.98,0.82);
		\end{tikzpicture}
		\caption{校准曲线示意}
	\end{figure}
\end{note}
\subsubsection{多分类}
与二分类问题不同，多分类问题中不存在唯一的正类与负类，因此通常需要先分别考察模型在每个类别上的表现，再通过适当的平均方式得到整体评价指标。\par
对于任意类别$k\in\{1,2,\dots,K\}(K\geqslant3)$，可将类别$k$视为正类，将其余类别合并为负类，并按照二分类问题中的定义计算$\operatorname{TP}_k$、$\operatorname{FP}_k$、$\operatorname{FN}_k$和$\operatorname{TN}_k$，进而得到该类别的$\operatorname{Precision}_k$、$\operatorname{Recall}_k$和$\operatorname{F1}_k$。这种处理方式称为one-versus-rest。\par
分别报告每个类别的评价指标能够完整反映模型在各类别上的表现，但当类别数量较多或需要使用单一数值比较模型时，需要进一步对各类别的指标进行汇总。常用的汇总方式包括Macro平均、Weighted平均和Micro平均。\par
\begin{definition}
	设$m_k$表示类别$k$的某个评价指标，该指标的Macro平均定义为：
	\begin{equation*}
		m_{\mathrm{macro}}=\frac{1}{K}\sum_{k=1}^{K}m_k
	\end{equation*}
\end{definition}
\begin{definition}
	设$m_k$表示类别$k$的某个评价指标，类别$k$的真实样本数为$n_k$，该指标的Weighted平均定义为：
	\begin{equation*}
		m_{\mathrm{weighted}}=\sum_{k=1}^{K}\frac{n_k}{n}m_k
	\end{equation*}
\end{definition}
\begin{definition}
	Micro平均定义为：
	\begin{gather*}
		\operatorname{Precision}_{\mathrm{micro}}=\frac{\sum\limits_{k=1}^{K}\operatorname{TP}_k}{\sum\limits_{k=1}^{K}(\operatorname{TP}_k+\operatorname{FP}_k)} \\
		\operatorname{Recall}_{\mathrm{micro}}=\frac{\sum\limits_{k=1}^{K}\operatorname{TP}_k}{\sum\limits_{k=1}^{K}(\operatorname{TP}_k+\operatorname{FN}_k)} \\
		\operatorname{F1}_{\mathrm{micro}}=2\times\frac{\operatorname{Precision}_{\mathrm{micro}}\operatorname{Recall}_{\mathrm{micro}}}{\operatorname{Precision}_{\mathrm{micro}}+\operatorname{Recall}_{\mathrm{micro}}}
	\end{gather*}
\end{definition}
\begin{note}
	在每个样本只有一个真实标签且模型只预测一个类别的单标签多分类问题中，每个预测正确的样本产生一个真阳性，而每个预测错误的样本同时产生一个假阳性和一个假阴性。因此：
	\begin{equation*}
		\sum_{k=1}^{K}\operatorname{FP}_k=\sum_{k=1}^{K}\operatorname{FN}_k=n-\sum_{k=1}^{K}\operatorname{TP}_k
	\end{equation*}
	从而：
	\begin{equation*}
		\operatorname{Precision}_{\mathrm{micro}}=\operatorname{Recall}_{\mathrm{micro}}=\operatorname{F1}_{\mathrm{micro}}=\operatorname{Accuracy}
	\end{equation*}
	所以在多分类问题中，同时报告Accuracy和Micro通常不会提供额外信息。\par
	若希望所有类别具有相同的重要性，应重点考察Macro平均；若希望评价结果反映实际类别分布，可以使用Weighted平均；若希望从所有样本层面汇总分类结果，则可以使用Micro平均。
\end{note}

\subsection{回归任务}
设真实值为$\{y_i\}_{i=1}^{n}$，模型预测值为$\{\hat{y}_i\}_{i=1}^{n}$。
\begin{definition}
  称：
  \begin{equation*}
    \operatorname{MAE}=\frac{1}{n}\sum_{i=1}^n|y_i-\hat{y}_i|
  \end{equation*}
  为\gls{MAE}。
\end{definition}
\begin{note}
  MAE计算预测误差的绝对值并取平均，因而对所有误差赋予相同的线性权重。它具有易于解释、与原数据同量纲等优点，并对单个异常值不敏感。
\end{note}
\begin{definition}
  称：
  \begin{equation*}
    \operatorname{MAPE}=\frac{1}{n}\sum_{i=1}^n\left|\frac{y_i-\hat{y}_i}{y_i}\right|
  \end{equation*}
  为\gls{MAPE}。
\end{definition}
\begin{note}
  MAPE 将误差按照真实值比例归一化，适合对预测误差的相对大小进行衡量。
\end{note}